% 原书分章PDF第17—19页：第15.5—15.6节，完整结果/讨论、表15.5及图15.2。
\section{PROBIT 结果}

我们使用上述四种方法分析 PROBIT 数据：标准的、基于似然的 g 计算，序贯 g 计算，IPW，以及 LTMLE。两种 g 计算算法和 LTMLE 算法均考虑纵向计数结局。LTMLE 算法采用两种不同的实现：第一种用逻辑回归估计所有结局模型组成部分，以及处理和删失概率；第二种用超级学习器（super learner，SL）估计同样的组成部分（Polley 等，2011；van der Laan 等，2007）。为此，我们使用 R 程序包 \texttt{SuperLearner}（Polley 和 van der Laan，2013）。所选的学习器库包括主效应逻辑回归、广义可加模型（Hastie，2011）、均值估计、最近邻算法（Peters 和 Hothorn，2009；仅在对二元结局建模时使用）、多元自适应回归样条模型（Milborrow 等，2014），以及逐步 AIC 程序（Venables 和 Ripley，2002，中的 \texttt{stepAIC}）。除 g 计算外，所有方法的标准误均按第15.4.2节，利用考虑聚集性的 EIF 方差计算；g 计算方法的标准误则使用成对簇自助法估计（Cameron 等，2008）。

首先，我们考察出生后第一年发生住院的时间区间数的期望之差，比较母乳喂养持续3—6个月与1—2个月、超过9个月与3—6个月，以及超过9个月与1—2个月这三组情况。表15.5给出估计值、标准误和95\%置信区间。对于第一组比较，所有方法的点估计都提示，截至12个月时发生住院的时间区间期望数有所下降，不过基于似然的 g 计算和 IPW 得到的置信区间跨过零。g 计算似乎对不同的参数化较为敏感，因为两种版本的点估计差异很大，分别为 $-0.03$ 和 $-0.12$。

对于第二组比较，除 IPW 以外，所有估计量的点估计均为零，提示将母乳喂养延长至3—6个月以后，并不会在住院方面带来更多益处。IPW 却估计，将母乳喂养从3—6个月延长至超过9个月，会增加住院。我们认为这一结论不合理。对于第三组比较，除 IPW 外，其他方法得到的住院减少幅度，基本与第一组比较相同，这也对应于：将母乳喂养持续时间延长到3—6个月之后，并未带来进一步改善。对于全部三组比较，两种 LTMLE 实现均给出相同的估计值和置信区间，提示超级学习器在本应用中没有改进估计。不过，对感染次数建模时，并非如此（Schnitzer 等，2014）。随着固定母乳喂养方案的持续时间增加，可以从图中观察到所估计的住院区间期望数如何变化。图15.2给出每种母乳喂养持续时间下，使用超级学习器的 LTMLE 所得边际均值估计及逐点置信区间。

\begin{table}[htbp]
\centering
\caption{不同母乳喂养持续时间下住院次数的边际期望}
\label{tab:ch15-probit-results}
\small
\begin{tabular}{lccc}
\toprule
方法 & 估计值 & 标准误 & 95\%置信区间\\
\midrule
\multicolumn{4}{c}{3—6个月与1—2个月比较}\\
g 计算（似然） & $-0.03$ & 0.02 & $(-0.06,0.00)$\\
g 计算（序贯） & $-0.12$ & 0.03 & $(-0.17,-0.06)$\\
IPW & $-0.05$ & 0.03 & $(-0.10,0.01)$\\
参数 LTMLE & $-0.06$ & 0.01 & $(-0.08,-0.04)$\\
使用 SL 的 LTMLE & $-0.06$ & 0.01 & $(-0.08,-0.04)$\\
\addlinespace
\multicolumn{4}{c}{9+个月与3—6个月比较}\\
g 计算（似然） & $-0.00$ & 0.01 & $(-0.01,0.01)$\\
g 计算（序贯） & $-0.00$ & 0.01 & $(-0.03,0.02)$\\
IPW & $0.05$ & 0.02 & $(0.01,0.08)$\\
参数 LTMLE & $-0.00$ & 0.01 & $(-0.02,0.02)$\\
使用 SL 的 LTMLE & $-0.00$ & 0.01 & $(-0.02,0.02)$\\
\addlinespace
\multicolumn{4}{c}{9+个月与1—2个月比较}\\
g 计算（似然） & $-0.03$ & 0.02 & $(-0.07,0.01)$\\
g 计算（序贯） & $-0.12$ & 0.03 & $(-0.17,-0.07)$\\
IPW & $0.00$ & 0.03 & $(-0.05,0.06)$\\
参数 LTMLE & $-0.06$ & 0.01 & $(-0.08,-0.04)$\\
使用 SL 的 LTMLE & $-0.06$ & 0.01 & $(-0.08,-0.04)$\\
\bottomrule
\end{tabular}
\end{table}

\begin{figure}[htbp]
\centering
% 图15.2：32条原始矢量绘制对象、6个估计点及两条置信界原样保留；只替换两个轴标签。
\begin{tikzpicture}[x=.53mm,y=-.53mm]
 \node[anchor=north west,inner sep=0pt] at (0,0)
  {\includegraphics[width=121.37mm,height=109.71mm]{figures/ch15_hospitalization_data.pdf}};
 \node[font=\small] at (127.7555,198.4015) {停止母乳喂养的时间};
 \node[font=\small,rotate=90] at (9.1171,90.1899) {住院次数的边际期望};
\end{tikzpicture}

\caption{使用 LTMLE 与超级学习估计的发生住院的时间区间期望数，包含95\%置信区域。}
\label{fig:ch15-hospitalization}
\end{figure}
\FloatBarrier

\section{讨论}

本章展示了如何调整 LTMLE 算法，使其适用于实际应用。具体而言，我们说明了如何修改 LTMLE，以更好地纳入纵向计数结构，以及如何修改基于 EIF 的方差估计，以适应多中心研究中的聚集性。对 PROBIT 数据的 LTMLE 分析显示，母乳喂养很可能影响婴儿出生后前12个月的住院。不过，我们估计，将母乳喂养延长至6个月以上，并不能获得额外益处。相较之下，我们此前的研究提示，将母乳喂养延长至9个月以后，12个月时胃肠道感染次数的边际期望仍然会继续下降（Schnitzer 等，2014）。在统计估计量的可靠性方面，我们发现 g 计算对参数化的选择较为敏感。IPW 估计量得出的结论是，将母乳喂养延长至6个月以后，会增加发生住院的时间区间期望数。这一结论缺乏科学合理性，也与其他估计量的结果矛盾。由于非稳定化权重的最大值为225，而第六时点有15,642名未删失研究对象，我们认为，IPW 的异常结果应当源于处理模型设定错误所造成的偏差，而非不稳定性。因此，LTMLE 能够利用结局模型组成部分修正这一估计，体现了采用这种更为复杂方法的益处。
\clearpage
\endgroup
